% ============================================================
% Notes: Fractional oscillator + Fractional wave equations
% Topics: fractional harmonic oscillator, intrinsic damping,
%         fractional wave equation, relaxation/localization, properties
% Source: Herrmann, Fractional Calculus: An Introduction for Physicists
% ============================================================

\section{Fractional oscillator model}

\subsection{Classical harmonic oscillator}
The classical oscillator is
\begin{equation}
\ddot{x}(t) + \omega^2 x(t)=0,
\end{equation}
with solutions $x(t)=A\cos(\omega t)+B\sin(\omega t)$.

\subsection{Fractional harmonic oscillator (general idea)}
A fractional oscillator replaces integer derivatives by fractional ones.
A standard prototype is
\begin{equation}
\boxed{
{}^{C}D_t^{2\alpha}x(t)+\omega^{2}x(t)=0,
\qquad 0<\alpha\le 1,
}
\end{equation}
where ${}^{C}D_t^{2\alpha}$ is the Caputo fractional derivative.
For $\alpha=1$ this reduces to the classical oscillator.

\subsection{Caputo fractional trigonometric functions (Herrmann)}
Herrmann constructs two linearly independent solutions of the fractional oscillator
and denotes them in analogy to sine/cosine:
\begin{align}
\boxed{
C\sin(\alpha,t)
=
\sum_{n=0}^{\infty}
(-1)^n
\frac{t^{(2n+1)\alpha}}{\Gamma(1+(2n+1)\alpha)},
}
\\[4pt]
\boxed{
C\cos(\alpha,t)
=
\sum_{n=0}^{\infty}
(-1)^n
\frac{t^{2n\alpha}}{\Gamma(1+2n\alpha)}.
}
\end{align}
These satisfy the key derivative relations (Caputo derivative):
\begin{align}
\boxed{
{}^{C}D_t^\alpha\,C\sin(\alpha,\omega t)=\omega^\alpha\,C\cos(\alpha,\omega t),
}
\\[4pt]
\boxed{
{}^{C}D_t^\alpha\,C\cos(\alpha,\omega t)=-\omega^\alpha\,C\sin(\alpha,\omega t).
}
\end{align}
Herrmann emphasizes these as the natural fractional analogues of trig functions. :contentReference[oaicite:4]{index=4}

\subsection{Mittag--Leffler representation}
The above series are special cases of Mittag--Leffler functions:
\begin{equation}
C\cos(\alpha,t)=E(2\alpha,-t^{2\alpha}),
\qquad
C\sin(\alpha,t)=t^\alpha\,E(2\alpha,1+\alpha,-t^{2\alpha}).
\end{equation}
Mittag--Leffler functions play the role of exponentials in fractional dynamics. :contentReference[oaicite:5]{index=5}

\subsection{Intrinsic damping (relaxation of oscillations)}
A crucial physical observation in Herrmann:
\begin{quote}
For $\alpha<1$, fractional oscillator solutions behave like \emph{damped oscillations}
even without adding an external friction term.
\end{quote}
So the damping is \textbf{intrinsic} to the fractional time evolution, not due to external forces. :contentReference[oaicite:6]{index=6} :contentReference[oaicite:7]{index=7}

\subsection{Definition-independence (robust qualitative behavior)}
Herrmann compares solutions for different fractional derivative definitions
and notes:
\begin{itemize}
\item for $\alpha\lesssim 1$, oscillations have similar frequency across definitions,
\item the main differences appear in the number of zeros (important for boundary problems).
\end{itemize}
:contentReference[oaicite:8]{index=8}

% ------------------------------------------------------------

\section{Fractional wave equations}

\subsection{Classical wave equation}
In 3D:
\begin{equation}
\left(\Delta - \frac{1}{c^2}\frac{\partial^2}{\partial t^2}\right)\psi(\mathbf{x},t)=0.
\end{equation}

\subsection{Fractional extension (spatial fractionalization)}
Herrmann proposes a fractional wave equation by replacing spatial second derivatives
with fractional ones (time kept second order):
\begin{equation}
\boxed{
\left(
\frac{d^{2\alpha}}{dx^{2\alpha}}
+
\frac{d^{2\alpha}}{dy^{2\alpha}}
+
\frac{d^{2\alpha}}{dz^{2\alpha}}
-\mu^2\frac{d^2}{dt^2}
\right)\psi(x,y,z,t)=0.
}
\end{equation}
Here $\mu$ carries dimensions $[s/m^\alpha]$. :contentReference[oaicite:9]{index=9}

\subsection{Separation of variables}
Using the product ansatz
\begin{equation}
\psi(x,y,z,t)=X(x)Y(y)Z(z)T(t),
\end{equation}
the PDE reduces to uncoupled ODEs:
\begin{align}
\boxed{
\frac{d^{2\alpha}}{dx^{2\alpha}}X = -\kappa_x^2 X,
\qquad
\frac{d^{2\alpha}}{dy^{2\alpha}}Y = -\kappa_y^2 Y,
\qquad
\frac{d^{2\alpha}}{dz^{2\alpha}}Z = -\kappa_z^2 Z,
}
\\[4pt]
\boxed{
\frac{d^2}{dt^2}T = \frac{\omega^2}{\mu^2}T.
}
\end{align}
The separation constants satisfy:
\begin{equation}
\boxed{
\omega=\mu\sqrt{\kappa_x^2+\kappa_y^2+\kappa_z^2}.
}
\end{equation}
:contentReference[oaicite:10]{index=10} :contentReference[oaicite:11]{index=11}

\subsection{Boundary conditions viewpoint}
Wave problems are often determined by boundary conditions (Dirichlet/Neumann/mixed),
not initial conditions.
Herrmann highlights that this makes Riemann-type solutions (even if singular at $t=0$)
still useful in wave contexts. :contentReference[oaicite:12]{index=12}

% ------------------------------------------------------------

\section{Parity, regularization, and localization}

\subsection{Why parity becomes nontrivial}
Fractional powers like $x^\alpha$ are multivalued for $x<0$:
\[
(-x)^\alpha = (-1)^\alpha |x|^\alpha
\]
which causes issues for defining physical wave solutions on domains containing $x<0$. :contentReference[oaicite:13]{index=13}

\subsection{Herrmann's parity-regularized coordinate}
To preserve parity behavior, Herrmann introduces a regularized coordinate
\begin{equation}
\boxed{
\hat{x}^\alpha := \mathrm{sign}(x)|x|^\alpha,
}
\end{equation}
and replaces
\begin{equation}
\boxed{
\frac{d}{dx}\ \longrightarrow\ \frac{d^\alpha}{d\hat{x}^\alpha}=: \hat{D}_x^\alpha.
}
\end{equation}
This yields a \emph{regularized fractional wave equation} with proper parity properties:
\begin{equation}
\boxed{
\left(
(\hat{D}_x^\alpha)^2+(\hat{D}_y^\alpha)^2+(\hat{D}_z^\alpha)^2
-\mu^2\frac{d^2}{dt^2}
\right)\psi(x,y,z,t)=0.
}
\end{equation}
Free solutions then transform under parity similarly to the classical ($\alpha=1$) case. :contentReference[oaicite:14]{index=14}

\subsection{Localization (qualitative meaning)}
In fractional wave/quantum problems, solutions may exhibit stronger localization due to:
\begin{itemize}
\item nonlocal spatial operators (long-range coupling),
\item altered nodal structure (different zero distributions),
\item modified dispersion relations (effective $k^{2\alpha}$ scaling).
\end{itemize}
In Herrmann's framework, the number/location of zeros is a major qualitative difference
between fractional and classical oscillatory solutions. :contentReference[oaicite:15]{index=15}

% ------------------------------------------------------------

\section{Relaxation vs localization}

\subsection{Relaxation (temporal damping)}
Fractional oscillators with $\alpha<1$ show \textbf{relaxation}:
\begin{itemize}
\item oscillations decay even without explicit friction,
\item decay is governed by Mittag--Leffler-type behavior,
\item memory effects are intrinsic to fractional derivatives.
\end{itemize}
This is one of the central physical signatures of fractional time dynamics. :contentReference[oaicite:16]{index=16}

\subsection{Localization (spatial fractionalization)}
Spatial fractional wave equations modify eigenmodes.
Under boundary conditions, the spatial fractional eigenfunctions can:
\begin{itemize}
\item deviate from ordinary sines/cosines,
\item show different node counts and spacing,
\item depend more sensitively on parity/regularization choices.
\end{itemize}
These effects can be interpreted as a form of localization/shape deformation
compared to classical wave modes. :contentReference[oaicite:17]{index=17} :contentReference[oaicite:18]{index=18}

% ------------------------------------------------------------

\section{Main properties (high-yield)}

\begin{itemize}
\item Fractional oscillator solutions are expressed via Mittag--Leffler functions
(the fractional analogue of exponentials). :contentReference[oaicite:19]{index=19}
\item For $\alpha<1$, the fractional harmonic oscillator behaves like a \textbf{damped oscillator}
with \textbf{intrinsic absorption} (no external friction needed). :contentReference[oaicite:20]{index=20}
\item Fractional wave equations can be formed by fractionalizing spatial derivatives,
keeping time second order. :contentReference[oaicite:21]{index=21}
\item Separation of variables still works and yields fractional eigenvalue ODEs. :contentReference[oaicite:22]{index=22}
\item Parity and negative arguments require care; Herrmann introduces a regularization
$\hat{x}^\alpha=\mathrm{sign}(x)|x|^\alpha$ to preserve parity symmetry. :contentReference[oaicite:23]{index=23}
\item Boundary conditions are central for wave problems; fractional solutions may differ
strongly in zero structure (important for mode spectra). :contentReference[oaicite:24]{index=24}
\end{itemize}
